\section*{NeurIPS Paper Checklist}

\begin{enumerate}

\item {\bf Claims}
    \item[] Question: Do the main claims made in the abstract and introduction accurately reflect the paper's contributions and scope?
    \item[] Answer: \answerYes{}
    \item[] Justification: The abstract and introduction state the paper's main contributions: CIR, CVN, the CIR-to-CVN translation, and the goal-aware generate-verify-repair loop. The paper also states the main scope restriction that the verifier reasons about validated, bounded CIR artifacts rather than arbitrary source code; see Sections~\ref{sec:introduction}, \ref{sec:method}, and the limitations discussion.

\item {\bf Limitations}
    \item[] Question: Does the paper discuss the limitations of the work performed by the authors?
    \item[] Answer: \answerYes{}
    \item[] Justification: The paper discusses that the current trust boundary is the generated CIR artifact, not generated source code; that the analysis assumes bounded thread instances and finite abstract data domains; and that scheduler-dependent properties such as livelock and starvation are outside the current guarantee. The evaluation is also described as a bounded, pattern-oriented benchmark rather than a corpus-scale empirical study.

\item {\bf Theory assumptions and proofs}
    \item[] Question: For each theoretical result, does the paper provide the full set of assumptions and a complete (and correct) proof?
    \item[] Answer: \answerYes{}
    \item[] Justification: The paper states the formal definitions of CIR, CVN, state, firing semantics, guard evaluation, goal satisfaction, and bug predicates in Sections~\ref{sec:method} and the appendices. The translation-correspondence results and their assumptions, including finite-domain boundedness and the treatment of unknown guards, are stated and proved in the appendix.

\item {\bf Experimental result reproducibility}
    \item[] Question: Does the paper fully disclose all the information needed to reproduce the main experimental results of the paper to the extent that it affects the main claims and/or conclusions of the paper (regardless of whether the code and data are provided or not)?
    \item[] Answer: \answerYes{}
    \item[] Justification: Section~\ref{sec:evaluation} specifies the benchmark patterns, evaluated LLMs, temperature, output-token limit, repair budget, acceptance criteria, hardware, and wall-clock measurement protocol. The appendix/supplementary material provides the CIR schemas, diagnostic format, and prompt structure needed to reproduce the reported generation and repair loop, subject to the availability and version stability of the closed-source LLM APIs.

\item {\bf Open access to data and code}
    \item[] Question: Does the paper provide open access to the data and code, with sufficient instructions to faithfully reproduce the main experimental results, as described in supplemental material?
    \item[] Answer: \answerYes{}
    \item[] Justification: The anonymized supplemental material includes the verifier implementation, benchmark CIR artifacts, benchmark specifications, prompt templates, diagnostic templates, and scripts for reproducing the reported translation, analysis, and repair-loop results. The supplement also documents the software environment and commands needed to run the verifier experiments.

\item {\bf Experimental setting/details}
    \item[] Question: Does the paper specify all the training and test details (e.g., data splits, hyperparameters, how they were chosen, type of optimizer) necessary to understand the results?
    \item[] Answer: \answerYes{}
    \item[] Justification: The paper does not train a model; the experiments evaluate a verification-and-repair pipeline using fixed LLM settings. Section~\ref{subsec:setup} reports the evaluated models, temperature, token limit, generation and repair budgets, benchmark construction, verifier acceptance criteria, and hardware environment.

\item {\bf Experiment statistical significance}
    \item[] Question: Does the paper report error bars suitably and correctly defined or other appropriate information about the statistical significance of the experiments?
    \item[] Answer: \answerNo{}
    \item[] Justification: The experiments report pattern-level success/failure counts, repair rounds, regressions, and wall-clock averages over repeated verifier runs, but do not report confidence intervals or statistical significance tests. This choice reflects the concept-and-feasibility nature of the benchmark, the small number of patterns, and the deterministic temperature-zero LLM setting; the paper does not claim statistically significant population-level superiority.

\item {\bf Experiments compute resources}
    \item[] Question: For each experiment, does the paper provide sufficient information on the computer resources (type of compute workers, memory, time of execution) needed to reproduce the experiments?
    \item[] Answer: \answerYes{}
    \item[] Justification: Section~\ref{subsec:setup} reports that the experiments were run on an Apple M4 Pro machine with 24GB memory, and the paper reports wall-clock averages for verifier execution. The formal verifier is CPU-only; LLM generation and repair calls are performed through hosted model APIs with a fixed round budget.

\item {\bf Code of ethics}
    \item[] Question: Does the research conducted in the paper conform, in every respect, with the NeurIPS Code of Ethics \url{https://neurips.cc/public/EthicsGuidelines}?
    \item[] Answer: \answerYes{}
    \item[] Justification: The research uses synthetic concurrency specifications and benchmark artifacts, does not involve human subjects, private personal data, scraped user data, or deployment in a real-world decision-making system. The paper preserves the trust boundary of the verifier and avoids claiming source-code correctness beyond the modeled CIR artifact.

\item {\bf Broader impacts}
    \item[] Question: Does the paper discuss both potential positive societal impacts and negative societal impacts of the work performed?
    \item[] Answer: \answerYes{}
    \item[] Justification: The paper discusses the positive potential of improving the reliability of LLM-assisted concurrent program construction through formal verification. It also discusses the main risk: users may over-trust a verified CIR artifact as evidence of source-code correctness, so the paper explicitly limits the guarantee to bounded CIR models and identifies source-code conformance checking as future work.

\item {\bf Safeguards}
    \item[] Question: Does the paper describe safeguards that have been put in place for responsible release of data or models that have a high risk for misuse (e.g., pre-trained language models, image generators, or scraped datasets)?
    \item[] Answer: \answerNA{}
    \item[] Justification: The paper does not release a high-risk pretrained model, image generator, scraped dataset, or user-derived dataset. The artifacts discussed are synthetic benchmark specifications, verification models, and verifier code; the evaluated LLMs are accessed through their existing providers and are not released by this work.

\item {\bf Licenses for existing assets}
    \item[] Question: Are the creators or original owners of assets (e.g., code, data, models), used in the paper, properly credited and are the license and terms of use explicitly mentioned and properly respected?
    \item[] Answer: \answerYes{}
    \item[] Justification: The paper identifies the evaluated LLMs and uses them through their respective providers' access mechanisms. No external dataset is used; the benchmark patterns and CIR artifacts are constructed by the authors. Existing research systems and libraries discussed in the paper are cited in the related work and experimental setup.

\item {\bf New assets}
    \item[] Question: Are new assets introduced in the paper well documented and is the documentation provided alongside the assets?
    \item[] Answer: \answerYes{}
    \item[] Justification: The paper introduces a small synthetic benchmark suite of bounded-concurrency patterns and associated CIR artifacts. These assets are documented in the paper and appendix, and the anonymized supplemental package includes the benchmark files, schema documentation, and reproduction instructions.

\item {\bf Crowdsourcing and research with human subjects}
    \item[] Question: For crowdsourcing experiments and research with human subjects, does the paper include the full text of instructions given to participants and screenshots, if applicable, as well as details about compensation (if any)? 
    \item[] Answer: \answerNA{}
    \item[] Justification: The paper does not involve crowdsourcing, human-subject experiments, user studies, surveys, or paid annotation work.

\item {\bf Institutional review board (IRB) approvals or equivalent for research with human subjects}
    \item[] Question: Does the paper describe potential risks incurred by study participants, whether such risks were disclosed to the subjects, and whether Institutional Review Board (IRB) approvals (or an equivalent approval/review based on the requirements of your country or institution) were obtained?
    \item[] Answer: \answerNA{}
    \item[] Justification: The paper does not involve human subjects, crowdsourcing participants, personal data collection, or behavioral experiments, so IRB or equivalent review is not applicable.

\item {\bf Declaration of LLM usage}
    \item[] Question: Does the paper describe the usage of LLMs if it is an important, original, or non-standard component of the core methods in this research? Note that if the LLM is used only for writing, editing, or formatting purposes and does \emph{not} impact the core methodology, scientific rigor, or originality of the research, declaration is not required.
    \item[] Answer: \answerYes{}
    \item[] Justification: LLMs are a core component of the proposed method: they generate CIR artifacts and repair them using structured diagnostics from the verifier. Section~\ref{subsec:setup} reports the evaluated models and inference settings, and Sections~\ref{subsec:pipeline} and \ref{subsec:repair-loop} describe how LLM generation and repair are integrated into the verification loop.

\end{enumerate}
